% ----------------------------------------------------
% Subsystem Design
% ----------------------------------------------------
\documentclass[class=report,11pt,crop=false]{standalone}
% Page geometry
\usepackage[a4paper,margin=20mm,top=25mm,bottom=25mm]{geometry}

\newcommand{\tabitem}{~~\llap{\textbullet}~~}
% Font choice
\usepackage{lmodern}

\usepackage{lipsum}

% Use IEEE bibliography style
\bibliographystyle{IEEEtran}

% Line spacing
\usepackage{setspace}
\setstretch{1.20}

% Ensure UTF8 encoding
\usepackage[utf8]{inputenc}

% Language standard (not too important)
\usepackage[english]{babel}

% Skip a line in between paragraphs
\usepackage{parskip}

% For the creation of dummy text
\usepackage{blindtext}

% Math
\usepackage{amsmath}

% Header & Footer stuff
\usepackage{fancyhdr}
\pagestyle{fancy}
\fancyhead{}
\fancyhead[R]{\nouppercase{\rightmark}}
\fancyfoot{}
\fancyfoot[C]{\thepage}
\renewcommand{\headrulewidth}{0.0pt}
\renewcommand{\footrulewidth}{0.0pt}
\setlength{\headheight}{13.6pt}

% Epigraphs
\usepackage{epigraph}
\setlength\epigraphrule{0pt}
\setlength{\epigraphwidth}{0.65\textwidth}

% Colour
\usepackage{color}
\usepackage[usenames,dvipsnames]{xcolor}

% Hyperlinks & References
\usepackage{hyperref}
\definecolor{linkColour}{RGB}{77,71,179}
\hypersetup{
    colorlinks=true,
    linkcolor=linkColour,
    filecolor=linkColour,
    urlcolor=linkColour,
    citecolor=linkColour,
}
\urlstyle{same}

% Automatically correct front-side quotes
\usepackage[autostyle=false, style=ukenglish]{csquotes}
\MakeOuterQuote{"}

% Graphics
\usepackage{graphicx}
\graphicspath{{Images/}{../Images/}}
\usepackage{makecell}
\usepackage{transparent}

% SI units
\usepackage{siunitx}

% Microtype goodness
\usepackage{microtype}

% Listings
\usepackage[T1]{fontenc}
\usepackage{listings}
\usepackage[scaled=0.8]{DejaVuSansMono}

% Custom colours for listings
\definecolor{backgroundColour}{RGB}{250,250,250}
\definecolor{commentColour}{RGB}{73, 175, 102}
\definecolor{identifierColour}{RGB}{196, 19, 66}
\definecolor{stringColour}{RGB}{252, 156, 30}
\definecolor{keywordColour}{RGB}{50, 38, 224}
\definecolor{lineNumbersColour}{RGB}{127,127,127}
\lstset{
  language=Matlab,
  captionpos=b,
  aboveskip=15pt,belowskip=10pt,
  backgroundcolor=\color{backgroundColour},
  basicstyle=\ttfamily,%\footnotesize,        % the size of the fonts that are used for the code
  breakatwhitespace=false,         % sets if automatic breaks should only happen at whitespace
  breaklines=true,                 % sets automatic line breaking
  postbreak=\mbox{\textcolor{red}{$\hookrightarrow$}\space},
  commentstyle=\color{commentColour},    % comment style
  identifierstyle=\color{identifierColour},
  stringstyle=\color{stringColour},
   keywordstyle=\color{keywordColour},       % keyword style
  %escapeinside={\%*}{*)},          % if you want to add LaTeX within your code
  extendedchars=true,              % lets you use non-ASCII characters; for 8-bits encodings only, does not work with UTF-8
  frame=single,	                   % adds a frame around the code
  keepspaces=true,                 % keeps spaces in text, useful for keeping indentation of code (possibly needs columns=flexible)
  morekeywords={*,...},            % if you want to add more keywords to the set
  numbers=left,                    % where to put the line-numbers; possible values are (none, left, right)
  numbersep=5pt,                   % how far the line-numbers are from the code
  numberstyle=\tiny\color{lineNumbersColour}, % the style that is used for the line-numbers
  rulecolor=\color{black},         % if not set, the frame-color may be changed on line-breaks within not-black text (e.g. comments (green here))
  showspaces=false,                % show spaces everywhere adding particular underscores; it overrides 'showstringspaces'
  showstringspaces=false,          % underline spaces within strings only
  showtabs=false,                  % show tabs within strings adding particular underscores
  stepnumber=1,                    % the step between two line-numbers. If it's 1, each line will be numbered
  tabsize=2,	                   % sets default tabsize to 2 spaces
  %title=\lstname                   % show the filename of files included with \lstinputlisting; also try caption instead of title
}

% Caption stuff
\usepackage[hypcap=true, justification=centering]{caption}
\usepackage{subcaption}

% Glossary package
% \usepackage[acronym]{glossaries}
\usepackage{glossaries-extra}
\setabbreviationstyle[acronym]{long-short}

% For Proofs & Theorems
\usepackage{amsthm}

% Maths symbols
\usepackage{amssymb}
\usepackage{mathrsfs}
\usepackage{mathtools}

% For algorithms
\usepackage[]{algorithm2e}

% Spacing stuff
\setlength{\abovecaptionskip}{5pt plus 3pt minus 2pt}
\setlength{\belowcaptionskip}{5pt plus 3pt minus 2pt}
\setlength{\textfloatsep}{10pt plus 3pt minus 2pt}
\setlength{\intextsep}{15pt plus 3pt minus 2pt}

% For aligning footnotes at bottom of page, instead of hugging text
\usepackage[bottom]{footmisc}

% Add LoF, Bib, etc. to ToC
\usepackage[nottoc]{tocbibind}

% SI
\usepackage{siunitx}

% For removing some whitespace in Chapter headings etc
\usepackage{etoolbox}
\makeatletter
\patchcmd{\@makechapterhead}{\vspace*{50\p@}}{\vspace*{-10pt}}{}{}%
\patchcmd{\@makeschapterhead}{\vspace*{50\p@}}{\vspace*{-10pt}}{}{}%
\makeatother
\begin{document}
% ----------------------------------------------------
\chapter{Subsystem Design} \label{ch:design}
\vspace{-1cm}
% =====================================================
\section{Design Decisions}


\subsection{Final Design}
The following design...
% 1. Provide a description of your solution and motivate why this design was chosen.
% 2. Provide the final schematic, make sure to include:
%       - labels and component values
%       - descriptions/comments on different parts of the schematic
%       - the completed schematic page setup: title, date, revision number, author name and surname.
%       - power flags on all power connections (3V, 5V, Gnd, etc)
% 2. Provide the final PCB, make sure to include:
%       - front, back and 3D view 
\begin{figure}[h]
    \centering
    \includegraphics[width=0.5\linewidth]{EEE3088F_latex_template/Figures/signature.jpg}
    \caption{Schematic}
    \label{fig:schematic}
\end{figure}

\begin{figure}[h]
     \centering
    % First subfigure
    \begin{subfigure}[b]{0.3\linewidth}
        \includegraphics[width=\linewidth]{EEE3088F_latex_template/Figures/signature.jpg}
        \caption{Front PCB}
        \label{fig:PCB_front}
    \end{subfigure}
    \hfill % Space between the subfigures
    % Second subfigure
    \begin{subfigure}[b]{0.3\linewidth}
        \includegraphics[width=\linewidth]{EEE3088F_latex_template/Figures/signature.jpg}
        \caption{Back PCB}
        \label{fig:PCB_back}
    \end{subfigure}
    \hfill % Space between the subfigures
    % Third subfigure
    \begin{subfigure}[b]{0.3\linewidth}
        \includegraphics[width=\linewidth]{EEE3088F_latex_template/Figures/signature.jpg}
        \caption{3D PCB}
        \label{fig:PCB_3D}
    \end{subfigure}
    \caption{PCB}
    \label{fig:PCB}
\end{figure}
% =====================================================
\section{Failure Management}

\begin{table}[h]
    \centering
    \caption{CAPTION} \label{tab:failuremanagement}
    \begin{tabular}{c|l}
        \hline
        \textbf{Name} & \textbf{Description} \\
        \hline
          &  \\
          &  \\
         &  \\
         &  \\
         \hline
    \end{tabular}
\end{table}

% =====================================================
\section{System Integration and Interfacing}
To integrate the subsystem with the rest of the system ....
\\
\begin{table}[h]
  \begin{center}
    \caption{Interfacing specifications}
    \label{tab:Interfacing}
    \begin{tabular}{ >{\centering\arraybackslash}m{3cm}  m{5cm} m{7cm}}
      \hline
      \textbf{Interface} & \textbf{Description} & \textbf{Pins/Output} \\   
      \hline
      I001 & SparkFun 9DoF IMU Breakout board to STM for data transfer (SPI) & \tabitem MISO: Breakout MISO* to STM PB14
      \newline\indent\tabitem MOSI: Breakout MOSI* to STM PB15\newline\indent\tabitem SCLK: Breakout SCLK* to STM PB13\newline\indent\tabitem CS: Breakout CS* to STM PB12 (GPIO)\newline\indent\tabitem Power: Breakout 1V8-5V5* to STM 3V\newline\indent\tabitem GND: Breakout GND* to STM GND\\
      \hline
      I002 & STM to FTDI for data transfer to computer (UART) &\tabitem STM PA2(RX) to FTDI RXD \newline\indent\tabitem STM PA3(TX) to FTDI TXD \\
      \hline 
      I003 & FTDI to computer for data transfer from STM & \tabitem Assuming the FTDI is connected to a Micro B USB port (which can be used to connect to a USB cable to connect to a computer), output via cable  \\
      \hline
    \end{tabular}
  \end{center}
\end{table}

% ----------------------------------------------------
\ifstandalone
\bibliography{../Bibliography/References.bib}
\fi
\end{document}
% ----------------------------------------------------